\chapter{Uniform Mordell--Lang bounds}
\label{chap:main}

% archon:covers MR4276287UniformityInMordellLangForCurves.lean

This chapter derives the uniform bounds on rational points and torsion packets
from the geometric and height estimates established earlier.

\section{External inputs}

\begin{lemma}[Northcott's theorem]
  \label{lem:northcott}
  \lean{MR4276287UniformityInMordellLangForCurves.northcott}
  \uses{def:ambient_setup}
  \textit{Source: standard (see, e.g., Hindry--Silverman, Theorem B.2.3; possibly Mathlib.NumberTheory.Heights).}
  For any number field \(K\), any positive integer \(D\), and any bound \(B\),
  there are only finitely many \(\overline{\mathbb{Q}}\)-points of a projective variety
  that have degree \(\le D\) over \(K\) and Weil height \(\le B\).
\end{lemma}

% SOURCE QUOTE PROOF: "By Northcott's Theorem implies that $\tau(s)$ comes from a finite set in $\mathbb{A}_{g}(\IQbar)$ that depends only on $g,d,$ and $\ell$." (read from references/MR4276287-mordell-lang-curves-source/RatPt.tex)
\begin{proof}
  Standard finiteness result: the number of algebraic integers of bounded degree
  and size is finite by the product formula. See Hindry--Silverman, Theorem B.2.3.
  The relevant instance here is applied to the moduli point \(\tau(s)\) in
  \(\mathbb{A}_{g,\ell}(\IQbar)\) viewed as a projective variety.
\end{proof}

\begin{lemma}[Rémond's quantitative Mordell-Lang bound]
  \label{lem:remond}
  \lean{MR4276287UniformityInMordellLangForCurves.remond}
  \uses{def:ambient_setup}
  \textit{Source: Rémond \cite{DPvarabII}, page 643. No equivalent currently in Mathlib; project-specific stub.}
  Let \(C\) be a smooth curve of genus \(g \ge 2\) over a number field \(F\),
  embedded in its Jacobian. For a finitely generated subgroup \(\Gamma \subset \mathrm{Jac}(C)(F)\)
  of rank \(\rho\), the number of \(F\)-rational points on \(C\) that lie in \(\Gamma\)
  is bounded by a constant depending only on \(g\), \([F:\mathbb{Q}]\), and \(\rho\).
\end{lemma}

\begin{proof}
  This is the quantitative version of the Faltings--Vojta theorem due to Rémond.
  The bound is effective: it combines the Vojta--Mumford counting argument with an
  explicit control of the height threshold. The result is applied in the proof of
  \cref{thm:bd_rat_intro} to handle the finitely many \(\IQbar\)-isomorphism classes
  of curves of bounded moduli height.
\end{proof}

\begin{theorem}[Raynaud's theorem on Manin--Mumford]
  \label{thm:raynaud_mm}
  \lean{MR4276287UniformityInMordellLangForCurves.raynaudManinMumford}
  \uses{def:ambient_setup}
  \textit{Source: Raynaud \cite{Raynaud}, Manin--Mumford conjecture. No equivalent currently in Mathlib; project-specific stub.}
  Let \(C\) be a smooth curve of genus \(g \ge 2\) over \(\IQbar\), embedded in its
  Jacobian \(J = \mathrm{Jac}(C)\) via a base point \(P_0\). The torsion subgroup
  \(J(\IQbar)_{\mathrm{tors}}\) intersects \(C(\IQbar) - P_0\) in finitely many points.
\end{theorem}

\begin{proof}
  This is Raynaud's proof of the Manin--Mumford conjecture (1983). The finiteness
  of \((C - P_0) \cap J_{\mathrm{tors}}\) follows from the theory of formal groups
  and the structure of the Galois representation on torsion points. The result is
  applied in the proof of \cref{thm:bd_tor_intro_nf} to handle the finitely many
  curves of small moduli height.
\end{proof}

\section{Introduction}

\begin{theorem}[Uniformity in Mordell--Lang for rational points]
  \label{thm:bd_rat_intro}
  \lean{MR4276287UniformityInMordellLangForCurves.thmBdRatIntro}
  \uses{thm:bd_fin_rank,lem:moduli_height_comparison}
  % SOURCE: Introduction of the paper (read from references/MR4276287-mordell-lang-curves-source/intro.tex)
  % SOURCE QUOTE: "Let \(g \ge 2\) and \(d \ge 1\) be integers. Then there exists a constant \(c = c(g, d) \ge 1\) with the following property. If \(C\) is a smooth curve of genus \(g\) defined over a number field \(F\) with \([F:\IQ] \le d\), then \(\#C(F) \le c^{1+\rho}\)."
  \textit{Source: Introduction of the paper.}
  The finite-rank theorem supplies the exponential-in-rank estimate once the
  Jacobian height is large enough. The moduli-height comparison converts the
  bounded-degree number-field hypothesis into a uniform lower bound on the
  relevant moduli height, so the constants depend only on \(g\) and \(d\).
\end{theorem}

% SOURCE QUOTE PROOF: "To prove the theorem we may assume $C(F)\not=\emptyset$. So fix
% $P_0\in C(F)$. We consider the Abel--Jacobi embedding
% $C-P_0\subset \mathrm{Jac}(C)$ defined over $F$. Then $\# C(F)\le \#
% (C_{F'}(\IQbar)-P_0)\cap\Gamma = \#
% (\mathfrak{C}_s(\IQbar)-P_0)\cap \Gamma$. If
% $h_{\overline{\mathbb{A}_g},\cM}(\tau(s))\ge c_1$, the theorem
% follows from Proposition~\ref{prop:premazur}.
% Note that in this case, the constant $c$ in \eqref{eq:CKexpinrank} is independent of $d$.
% So we may assume that the height
% of $\tau(s)$ is less than $c_1$.
% As $[F':\IQ]\le [F':F][F:\IQ]\le [F':F]d$,
% Northcott's Theorem implies
% that $\tau(s)$ comes   from a finite set in $\mathbb{A}_{g}(\IQbar)$
% that depends only on $g,d,$ and $\ell$. The same holds for $s$ and
% thus $F'=F(s)$
% since the Torelli morphism $\tau$ is finite-to-$1$ and thus   has
% fibers of bounded cardinality.
% This means that the remaining $C$ are twists   in
% finitely many $F'$-isomorphism classes. But then it suffices to
% apply R\'emond's estimate \cite[page 643]{DPvarabII} to a single
% $C_{F'}$ and    use $\#C(F) \le \# (C_{F'}(\IQbar)-P_0)\cap\Gamma$ to
% conclude the theorem."
% (read from references/MR4276287-mordell-lang-curves-source/RatPt.tex)
\begin{proof}
  \uses{thm:bd_fin_rank,lem:moduli_height_comparison,lem:northcott,lem:remond}
  Fix \(P_0 \in C(F)\) and take \(\Gamma = \mathrm{Jac}(C)(F)\) with rank
  \(\rho\). The Abel--Jacobi embedding gives \(\#C(F) \le \#(\mathfrak{C}_s(\IQbar)
  - P_0) \cap \Gamma\).

  By \cref{lem:moduli_height_comparison}, the height \(h(\iota([\mathrm{Jac}(C)]))\)
  is comparable, up to affine-linear constants depending only on \(g\), to the
  height \(h_{\overline{\mathbb{A}_g},\cM}(\tau(s))\) on the fine moduli space.
  Since \([F:\mathbb{Q}] \le d\), the degree constraint places the moduli point
  in a region where \cref{thm:bd_fin_rank} applies: if
  \(h(\iota([\mathrm{Jac}(C)])) \ge c_1(g,\iota)\), then
  \(\#(C(\IQbar) - P_0) \cap \Gamma \le c_2^{1+\rho}\) with constants
  depending only on \(g\).

  If instead \(h(\iota([\mathrm{Jac}(C)])) < c_1\), Northcott's theorem and
  the finiteness of the Torelli fibers imply that \(\tau(s)\) comes from a
  finite set depending only on \(g\), \(d\), and the level. The remaining curves
  split into finitely many \(\IQbar\)-isomorphism classes, each handled
  individually by Rémond's quantitative bound. In both cases the final constant
  \(c = c(g,d)\) depends only on \(g\) and \(d\).
\end{proof}

\begin{theorem}[Finite-rank Mordell--Lang bound]
  \label{thm:bd_fin_rank}
  \lean{MR4276287UniformityInMordellLangForCurves.thmBdFinRank}
  \uses{prop:premazur,lem:moduli_height_comparison}
  % SOURCE: Introduction of the paper (read from references/MR4276287-mordell-lang-curves-source/intro.tex)
  % SOURCE QUOTE: "Let \(g \ge 2\) and let \(\iota\) be as above. Then there exist two constants \(c_1= c_1(g,\iota) \ge 0\) and \(c_2 = c_2(g,\iota) \ge 1\) with the following property."
  \textit{Source: Introduction of the paper.}
  The proof uses the pre-Mazur proposition on the specific curve obtained from
  the Jacobian at the moduli point, then compares the Jacobian height with the
  projective height of the moduli point to remove any dependence on the curve
  beyond the large-height threshold.
\end{theorem}

% SOURCE QUOTE PROOF: "Let $C$ be a smooth curve of genus $g\ge 2$ defined over $\IQbar$,
% and let $\Gamma$ be a finite rank subgroup of
% $\mathrm{Jac}(C)(\IQbar)$. Let $P_0 \in C(\IQbar)$.
% The curve $C$
% corresponds to a $\IQbar$-point $s_{\mathrm{c}}$ of
% $\mathbb{M}_{g,1}$.
% The fine moduli space $\mathbb{M}_g$ of smooth genus $g$ curves
% with level-$\ell$-structure is a finite covering of $\mathbb{M}_{g,1}$. So there exists an $s \in \mathbb{M}_g(\IQbar)$
% that maps to $s_{\mathrm{c}}$. Thus $C$ is isomorphic, over $\IQbar$, to the fiber $\mathfrak{C}_s$ of the universal family
% $\mathfrak{C}_g\rightarrow\mathbb{M}_g$. We thus view $\Gamma$ as a
% finite rank subgroup of $\mathrm{Jac}(\mathfrak{C}_s)(\IQbar)$, and
% $P_0 \in \mathfrak{C}_s(\IQbar)$.
% Consider the Abel--Jacobi embedding $C-P_0 \subseteq
% \mathrm{Jac}(C)$. Then $\# (C(\IQbar)-P_0) \cap \Gamma = \#
% (\mathfrak{C}_s(\IQbar) - P_0)\cap \Gamma$. If $h_{\overline{\mathbb{A}_g},\cM}(\tau(s)) \ge c_1$, then $\# (C(\IQbar)-P_0) \cap \Gamma \le c_2^{1+\rho}$
% by Proposition~\ref{prop:premazur}.
% Thus it suffices to find a constant $c'_1\ge 0$ that is independent
% of $C$ and such
% that $h(\iota([\mathrm{Jac}(C)])) \ge c_1'$ implies
% $h_{\overline{\mathbb{A}_g},\cM}(\tau(s)) \ge c_1$.
% As after the proof of Theorem~\ref{ThmBdRatIntro}
% denote by $\rho \colon \mathbb{A}_g = \mathbb{A}_{g,\ell}
% \rightarrow \mathbb{A}_{g,1}$ the natural morphism.
% As in the proof of Theorem~\ref{ThmBdFinRank} we use
% $h(\iota(\rho(t)))\le c'h_{\overline{\IA_{g}},\cM}(t)  + c''$ for all $t\in
% \mathbb{A}_g(\IQbar)$.
% The theorem follows since $\rho(\tau(s)) = [\mathrm{Jac}(C)]$."
% (read from references/MR4276287-mordell-lang-curves-source/RatPt.tex)
\begin{proof}
  \uses{prop:premazur,lem:moduli_height_comparison}
  Let \(C\) be a smooth curve of genus \(g \ge 2\) defined over \(\IQbar\),
  let \(\Gamma \subset \mathrm{Jac}(C)(\IQbar)\) have finite rank \(\rho\),
  and let \(P_0 \in C(\IQbar)\). Since \(\mathbb{M}_g\) is a finite cover of
  the coarse moduli space \(\mathbb{M}_{g,1}\), there exists
  \(s \in \mathbb{M}_g(\IQbar)\) with \(C \cong \mathfrak{C}_s\) over \(\IQbar\).
  We identify \(\Gamma\) with a finite-rank subgroup of
  \(\mathrm{Jac}(\mathfrak{C}_s)(\IQbar)\).

  Via the Abel--Jacobi embedding,
  \(\#(C(\IQbar) - P_0) \cap \Gamma = \#(\mathfrak{C}_s(\IQbar) - P_0) \cap \Gamma\).
  If \(h_{\overline{\mathbb{A}_g},\cM}(\tau(s)) \ge c_1\), then \cref{prop:premazur}
  gives \(\#(\mathfrak{C}_s(\IQbar) - P_0) \cap \Gamma \le c_2^{1+\rho}\)
  directly.

  It remains to convert the hypothesis \(h(\iota([\mathrm{Jac}(C)])) \ge c_1'\)
  into the moduli-height lower bound required by \cref{prop:premazur}. By
  \cref{lem:moduli_height_comparison}, the height \(h(\iota(\rho(\tau(s))))\)
  is bounded above by \(c' h_{\overline{\mathbb{A}_g},\cM}(\tau(s)) + c''\)
  for constants \(c' > 0\) and \(c'' \ge 0\) depending only on \(g\) and the
  immersions. Since \(\rho(\tau(s)) = [\mathrm{Jac}(C)]\), choosing
  \(c_1' = (c_1 + c'')/c'\) ensures that
  \(h(\iota([\mathrm{Jac}(C)])) \ge c_1'\) implies
  \(h_{\overline{\mathbb{A}_g},\cM}(\tau(s)) \ge c_1\), completing the proof.
\end{proof}

\begin{theorem}[Uniform torsion packets]
  \label{thm:bd_tor_intro_nf}
  \lean{MR4276287UniformityInMordellLangForCurves.thmBdTorIntroNF}
  \uses{thm:bd_fin_rank,lem:moduli_height_comparison}
  % SOURCE: Introduction of the paper (read from references/MR4276287-mordell-lang-curves-source/intro.tex)
  % SOURCE QUOTE: "Let \(g \ge 2\) and \(d \ge 1\) be integers. Then there exists a constant \(c = c(g, d) \ge 1\) with the following property. Let \(C\) be a smooth curve of genus \(g\) defined over a number field \(F\subset\IQbar\) with \([F:\IQ] \le d\) and let \(P_0 \in C(\IQbar)\), then \(\# (C(\overline \IQ)-P_0) \cap \mathrm{Jac}(C)(\IQbar)_{\mathrm{tors}} \le c\)."
  \textit{Source: Introduction of the paper.}
  This is the rank-zero specialization of the finite-rank theorem, combined with
  the same moduli-height comparison used for rational points. The base point
  disappears because torsion is a translate-invariant rank-zero subgroup.
\end{theorem}

% SOURCE QUOTE PROOF: "Apply Theorem~\ref{ThmBdFinRank} to $C_{\IQbar}$, $P_0 \in C(\IQbar)$
% and $\Gamma = \jac(C)({\IQbar})_{\mathrm{tors}}$, whose rank is $0$. Then we obtain $c_1 \ge 0$ and $c_2 \ge 1$ such that
%   \[
%   \#(C(\IQbar) - P_0) \cap \jac(C)({\IQbar})_{\mathrm{tors}} \le c_2
%   \]
%   if $h(\iota([\mathrm{Jac}(C_\IQbar)])) \ge c_1$.
%
% By the Northcott property and Torelli's Theorem, there are up-to
% $\IQbar$-isomophism only finitely many $C_{\IQbar}$'s  defined over a number field $F$ with $[F:\IQ] \le d$ such that  $h(\iota([\mathrm{Jac}(C_\IQbar)])) < c_1$. By applying Raynaud's result on the Manin--Mumford Conjecture to each one of these finitely many curves separately, we obtain Theorem~\ref{thm:BdTorIntroNF}."
% (read from references/MR4276287-mordell-lang-curves-source/RatPt.tex)
\begin{proof}
  \uses{thm:bd_fin_rank,lem:moduli_height_comparison,lem:northcott,thm:raynaud_mm}
  Apply \cref{thm:bd_fin_rank} with \(\Gamma = \mathrm{Jac}(C)(\IQbar)_{\mathrm{tors}}\),
  which has rank \(\rho = 0\). The theorem yields constants \(c_1 \ge 0\) and
  \(c_2 \ge 1\) such that, whenever \(h(\iota([\mathrm{Jac}(C)])) \ge c_1\),
  \[
    \#(C(\IQbar) - P_0) \cap \mathrm{Jac}(C)(\IQbar)_{\mathrm{tors}} \le c_2^{1+0} = c_2.
  \]
  By \cref{lem:moduli_height_comparison} and the degree bound \([F:\mathbb{Q}] \le d\),
  the constant \(c_2\) depends only on \(g\) and \(d\).

  If \(h(\iota([\mathrm{Jac}(C)])) < c_1\), Northcott's theorem and Torelli's
  theorem together imply that there are only finitely many
  \(\IQbar\)-isomorphism classes of curves \(C_{\IQbar}\) defined over a number
  field \(F\) with \([F:\mathbb{Q}] \le d\) satisfying this height bound. Applying
  Raynaud's theorem on the Manin--Mumford conjecture to each of these finitely
  many curves individually yields a uniform bound on the torsion packet,
  completing the proof.
\end{proof}

\section{Proof of Theorems~\ref{ThmBdRatIntro}, \ref{ThmBdFinRank}, and \ref{thm:BdTorIntroNF}}

\begin{proposition}[Pre-Mazur bound]
  \label{prop:premazur}
  \lean{MR4276287UniformityInMordellLangForCurves.preMazur}
  \uses{thm:ht_inequality_full,thm:nondeg_for_bd,lem:vojtamumford,prop:alg_pt_far}
  % SOURCE: Section 8 of the paper (read from references/MR4276287-mordell-lang-curves-source/RatPt.tex)
  % SOURCE QUOTE: "The exist constants \(c_1\ge 0,c_2\ge 1\) ... \(\# (\mathfrak{C}_s(\IQbar)-P_0)\cap \Gamma \le c_2^{1+\rho}\)."
  \textit{Source: Section 8 of the paper.}
  The proof fixes the moduli point, transfers the sparse-point statement to the
  corresponding curve in its Jacobian, and then splits the points into the large
  part controlled by Vojta--Mumford and the residual part controlled by the
  packet-counting and sparse-points lemmas.
\end{proposition}

% SOURCE QUOTE PROOF: "We now follow the argumentation in~\cite{DGH1p}. Let $\Gamma$ be a
% subgroup of $\mathfrak{A}_{g,\tau(s)}(\IQbar)$ for finite rank $\rho$.
% We first prove the proposition in the case $P_0=P_s$. We apply Lemma
% \ref{lem:vojtamumford}  to the curve $C =
% \mathfrak{C}_s-P_s\subset \mathfrak{A}_{g,\tau(s)}=A$ and use the bounds
% (\ref{eq:choosecNT}), \eqref{eq:h1bound}, and
% (\ref{eq:degreehgtbound}).
% Note that $C$ is a smooth curve of genus
% $g\ge 2$. So it cannot be the translate of an algebraic subgroup of
% $A$.
% It follows that the number of points $P\in
% \mathfrak{C}_s(\IQbar)$ with $P-P_s\in \Gamma$ and $\hat h(P-P_s) >
% R^2$ where
% \begin{equation}
%     R= (c\max\{1,h_{\overline{\mathbb{A}_g},\cM}(\tau(s))\})^{1/2}
% \end{equation}
% is at most  $c^{\rho}\le c^{1+\rho}$.
% The burden of this paper is  to find a bound of the same quality
% for the number of pairwise distinct points $P_1,P_2,P_3,\ldots$  in
% $\mathfrak{C}_s(\IQbar)$ with $\hat h(P_i-P_s) \le R^2$. This is where
% Proposition~\ref{PropAlgPtFar} enters. Recall our assumption
% (\ref{eq:hslbc1}) on $s$.
% Let $c_2'$ be the constant $c_2$ from Proposition~\ref{PropAlgPtFar}; it is independent of $s$. As
% $\# \Xi_s\le c_2'$ we may assume $P_i\not\in \Xi_s$ for all $i$.
% So we may assume that each $P_i$ is in the second alternative of
% Proposition~\ref{PropAlgPtFar}.
% As in \cite{DGH1p} we use the Euclidean norm $|\cdot|$ defined by $\hat h^{1/2}$ on
% the $\rho$-dimensional  $\IR$-vector space $\Gamma\otimes\IR$.
% Let $r\in (0,R]$. By an elementary ball packing argument, any subset of
% $\Gamma\otimes\IR$ contained in a closed ball of radius $R$ is
% covered by at most $(1+2R/r)^{\rho}$ closed balls of radius $r$ centered at
% elements of the given subset; see \cite[Lemme~6.1]{Remond:Decompte}.
% We apply this geometric argument to $R$ as in (\ref{def:bigR}) and
% to  $r$, the positive square-root of
% $h_{\overline{\mathbb{A}_g},\cM}(\tau(s))/c_3 =
% \max\{1,h_{\overline{\mathbb{A}_g},\cM}(\tau(s))\}/c_3$.
% The contribution of the height
% $h_{\overline{\mathbb{A}_g},\cM}(\tau(s))$ cancels in the quotient
% $R/r$. We find $R/r\le c$. So the
% number of balls in the covering is at most $c^{1+\rho}$.
% By Proposition~\ref{PropAlgPtFar}(ii) the number of the $P_i$'s that
% map to a single closed ball of radius $r$ is
% at most $c_4$. Thus after increasing $c$ we find that $\#\{P_i \in \mathfrak{C}_s(\IQbar)\cap \Gamma : \hat{h}(P_i-P_s) \le R^2\} \le c_4 c^{1+\rho}$,
% as desired.
% This completes the proof of the proposition in the case $P_0=P_s$
% for sufficiently large $c_2$.
% The case of a general base point follows easily as our estimates
% depend only on the rank $\rho$ of $\Gamma$. Indeed,
% let $P_0\in \mathfrak{C}_s(\IQbar)$ be an arbitrary point and let
% $\Gamma'$ be the subgroup of $\mathfrak{A}_{g,\tau(s)}(\IQbar)$ generated by
% $\Gamma$ and $P_0-P_s$. Its rank is at most $\rho +1$.
% Now if $Q\in \mathfrak{C}_s(\IQbar)-P_0$ lies in $\Gamma$, then
% $Q+P_0-P_s\in \mathfrak{C}_s(\IQbar)-P_s$ lies in $\Gamma'$. The
% number of such $Q$ is at most $c_2^{2+\rho}$ by
% what we already proved.
% The proposition follows as $c_2^{2+\rho}\le (c_2^2)^{1+\rho}$ and
% since we may replace $c_2$ by $c_2^2$."
% (read from references/MR4276287-mordell-lang-curves-source/RatPt.tex)
\begin{proof}
  \uses{thm:ht_inequality_full,thm:nondeg_for_bd,lem:vojtamumford,prop:alg_pt_far}
  The proof proceeds in two stages: bounding points of large Néron--Tate height
  and bounding points of small Néron--Tate height.

  \textbf{Auxiliary bounds.}
  Applying the Silverman--Tate comparison to the universal abelian family gives
  \(|{\hat{h}}(P) - h(P)| \le c\max\{1, h_{\overline{\mathbb{A}_g},\cM}(\tau(s))\}\)
  for all \(P \in \mathfrak{A}_{g,\tau(s)}(\IQbar)\), so we may take
  \(c_{\mathrm{NT}} = c\max\{1, h_{\overline{\mathbb{A}_g},\cM}(\tau(s))\}\).
  A torsion-point argument using Cramer's rule produces the addition-law data
  \((P_{ij})\) with height \(h_1 \le c\max\{1, h_{\overline{\mathbb{A}_g},\cM}(\tau(s))\}\).
  From \cref{thm:nondeg_for_bd} and the universal-family setup the translated
  curve satisfies \(\deg(\mathfrak{C}_s - P_s) \le c\) and
  \(h(\mathfrak{C}_s - P_s) \le c\max\{1, h_{\overline{\mathbb{A}_g},\cM}(\tau(s))\}\).

  \textbf{Large-height points.}
  Apply \cref{lem:vojtamumford} to \(C = \mathfrak{C}_s - P_s \subset
  \mathfrak{A}_{g,\tau(s)}\) with the bounds for \(c_{\mathrm{NT}}\), \(h_1\),
  and \(\deg C\). Since \(C\) is a smooth curve of genus \(g \ge 2\), it cannot
  be a translate of an algebraic subgroup, so \cref{lem:vojtamumford} applies.
  Setting \(R = (c\max\{1, h_{\overline{\mathbb{A}_g},\cM}(\tau(s))\})^{1/2}\),
  the number of points \(P \in \mathfrak{C}_s(\IQbar)\) with \(P - P_s \in \Gamma\)
  and \(\hat{h}(P - P_s) > R^2\) is at most \(c^\rho \le c^{1+\rho}\).

  \textbf{Small-height points.}
  For points \(P_i\) with \(\hat{h}(P_i - P_s) \le R^2\), apply
  \cref{prop:alg_pt_far} to \(S = \mathbb{M}_g\) using the height lower
  bound \(h_{\overline{\mathbb{A}_g},\cM}(\tau(s)) \ge c_1\) from
  \cref{thm:ht_inequality_full}. The exceptional set \(\Xi_s\) satisfies
  \(\#\Xi_s \le c_2'\) uniformly, so we may assume each \(P_i \notin \Xi_s\),
  placing every \(P_i\) in alternative~(ii) of \cref{prop:alg_pt_far}.
  Setting \(r = (h_{\overline{\mathbb{A}_g},\cM}(\tau(s))/c_3)^{1/2}\),
  any ball of Néron--Tate radius \(r\) in \(\Gamma \otimes \IR\) contains at
  most \(c_4\) of the \(P_i\). The ratio \(R/r \le c\) is bounded independently
  of \(s\) because the height cancels. A ball-packing argument covers the ball
  of radius \(R\) by at most \((1 + 2R/r)^\rho \le c^{1+\rho}\) balls of radius
  \(r\), so the total number of small-height points is at most \(c_4 c^{1+\rho}\).

  \textbf{General base point.}
  For an arbitrary \(P_0 \in \mathfrak{C}_s(\IQbar)\), let \(\Gamma'\) be
  generated by \(\Gamma\) and \(P_0 - P_s\); its rank is at most \(\rho + 1\).
  If \(Q \in \mathfrak{C}_s(\IQbar) - P_0\) lies in \(\Gamma\), then
  \(Q + P_0 - P_s \in \mathfrak{C}_s(\IQbar) - P_s\) lies in \(\Gamma'\).
  Applying the case \(P_0 = P_s\) with rank \(\rho + 1\) gives at most
  \(c_2^{2+\rho}\) such points, and \(c_2^{2+\rho} \le (c_2^2)^{1+\rho}\),
  so replacing \(c_2\) by \(c_2^2\) concludes the proof.
\end{proof}
